% SPDX-FileCopyrightText: Copyright (c) 2026 NVIDIA CORPORATION & AFFILIATES. All rights reserved.
% SPDX-License-Identifier: Apache-2.0
%
% ROLE (operator request R1, CR/facts/PAPER_REQUESTS.md): routing primitives and the Dynamo
%   worker-selection API; the default cost function in that vocabulary; how each pedagogical cost
%   function (the ported policies) fits under it; the learned policy as a generalization; and the
%   major caveats the campaign uncovered, each with evidence and status. Split for a clean main text
%   (restructure plan, base <commit-50>): the main text keeps the API's composition, the default cost,
%   tab:api-policies, the learned generalization and a highlight of the caveats that shape results;
%   eligibility, declared inputs, request facts, catalog loading, the default picker's details and
%   the per-port rules are in appendix-api.tex (app:api, app:ports), and the full caveat table is
%   in appendix-caveats.tex (app:caveats, tab:caveats).
% EVIDENCE: WT/lib/kv-router/src/plugins/worker_selection.rs and worker_selection/{context,inputs,
%   config}.rs; WT/lib/kv-router/src/scheduling/{filter.rs,selector/policy.rs,config.rs};
%   WT/lib/router-plugins/builtin/src/{default/,default.rs,two_tier_cost_fn.rs,lmetric.rs,ramjet.rs,
%   dualmap.rs,chwbl.rs,llm_d/,sticky_session.rs,thunderagent/,learned_choice/}; caveats:
%   CR/facts/{UPSTREAM_FOLLOWUPS.md,DEVIATIONS.md,setup.json,agentx_lowered.json,lag_patch.json,
%   robustness.json,live_smoke.json}, CR/audits/{setup,build,phase2-mid,final}/.
% STATUS: final for this draft.

\section{Background: routing primitives and Dynamo's worker-selection API}
\label{sec:background}

Every routing rule in this study, Dynamo's default, the ported heuristics and the learned policy,
runs through one interface: Dynamo's worker-selection API, a contract between the router
\emph{host}, which owns the request path and the shared state, and a \emph{policy}, which turns that
state into a choice. This section describes the interface as a small set of routing primitives,
writes the default cost function and every other catalog policy in its vocabulary, shows how the
learned policy generalizes them, and highlights the caveats that bound the results. The vocabulary
is the one the code uses.\footnote{Type and method names are those of the router's plugin API on
  the study's branch.}
% src: WT/lib/kv-router/src/plugins/worker_selection.rs (module doc: "Public filter, scorer, and
%      picker contracts and their input signals")

% ---------------------------------------------------------------------------------------------
\subsection{The decision and what the host owns}
\label{sec:api-host}

For each request $j$ the host builds a \emph{candidate table} of eligible workers; eligibility and
pins are the host's and always win over a policy's preference (\cref{app:api}).
% src: WT/lib/kv-router/src/scheduling/filter.rs (RoutingEligibility); scheduling/selector/policy.rs
%      ("Explicit request pins remain mandatory regardless of this option")
The host also owns the state a policy reads. It keeps the KV index that the engines' cache events
feed and computes each candidate's prefix overlap with the request, per cache tier; it books the
load of every request it dispatches (prefill tokens in flight, decode blocks, active requests); and
it knows each worker's configured KV capacity. A policy owns only its decision and whatever private
state it chooses to keep, such as a random generator, a session map or a sliding window. This split
is what makes a policy portable: the same policy file runs wherever a host can fill the inputs it
declares, including the offline replay used here (\cref{sec:aisim}).

% ---------------------------------------------------------------------------------------------
\subsection{Filters, scorers and pickers}
\label{sec:api-components}

A worker-selection policy is an ordered composition of three kinds of component
(\cref{fig:api}).
\begin{description}
  \item[\code{WorkerFilter}] keeps or drops one eligible candidate. Filters run in declaration
    order and can only shrink the set.
  \item[\code{WorkerScorer}] writes one finite, lower-is-better cost for every surviving candidate
    in a single call. The host validates each contribution and adds the contributions of all
    scorers, so scorers compose additively.
  \item[\code{WorkerPicker}] receives the surviving candidates with their total costs and returns
    one row. It may ignore the costs entirely.
\end{description}
% src: worker_selection.rs (WorkerFilter::keep, WorkerScorer::score "Write one finite,
%      lower-is-better contribution", WorkerPicker::pick); selector/policy.rs (ComposedPolicyState:
%      filters, scorers, picker; WorkerSelectionPolicy::new_with_filters)
None of the builtin policies in this study uses a filter; candidate restriction, where a policy
needs it, happens inside its picker. Every builtin policy is therefore ``scorers plus a picker'',
and the interesting differences are in which inputs each reads and whether its choice for one
worker depends on the other candidates.
% src: grep of WT/lib/router-plugins/builtin/src for "impl WorkerFilter for": no builtin filter

A component declares the optional worker signals it reads, such as cache overlap (\code{CACHE}),
router-tracked load (\code{LOAD}), a modeled prefill backlog (\code{PREFILL\_TIME}) and a cost
multiplier from the request's preferred routing constraints (\code{PREFERRED\_TAINT}), and the host
materializes only those; a policy also reads request facts such as the prompt length, its prefix
hashes and the session context (\cref{app:api}).
% src: worker_selection/inputs.rs (WorkerInputs); worker_selection/context.rs (WorkerSelectionContext)

\paragraph{The policy catalog.}
A deployment selects policies with a YAML file. Its \code{worker\_selection} block names one
instance per worker role (aggregated, prefill, decode, encode) and lists named instances, each a
catalog \code{type} with a \code{parameters} mapping (\cref{fig:api}); \cref{app:api} describes how
the host loads and validates the catalog.\footnote{The offline replay host on this study's branch
  runs the same catalog; on Dynamo's main branch at the time of the study, replay rejected catalog
  policies selected by YAML.}
% src (footnote): CR/facts/UPSTREAM_FOLLOWUPS.md "Straight-main lane outcomes" #1 ("on main replay
%      rejects YAML-selected catalog policies")
% src: worker_selection/config.rs (RawWorkerSelectionConfig: aggregated, prefill, decode, encode,
%      instances)

\begin{figure}[tbp]
  \centering
  \begin{tikzpicture}[
      node distance=4mm and 5mm,
      box/.style={draw, rounded corners=2pt, align=center, font=\scriptsize, inner sep=3pt,
        minimum height=8mm},
      host/.style={box, fill=black!6},
      arr/.style={-{Stealth[length=1.6mm]}, thick}]
    \node[host] (req) {request $j$\\prompt, prefix hashes,\\session, pins};
    \node[host, right=of req] (elig) {host eligibility\\allowed, available,\\overload, pins};
    \node[host, right=of elig] (table) {candidate table\\\code{CACHE} \code{LOAD}\\capacity, \dots};
    \node[box, right=of table] (filters) {filters\\keep / drop};
    \node[box, right=of filters] (scorers) {scorers\\$\sum$ costs};
    \node[box, right=of scorers] (picker) {picker\\one row};
    \draw[arr] (req) -- (elig);
    \draw[arr] (elig) -- (table);
    \draw[arr] (table) -- (filters);
    \draw[arr] (filters) -- (scorers);
    \draw[arr] (scorers) -- (picker);
    \node[font=\scriptsize, below=2mm of $(elig.south)!0.5!(table.south)$] {host-owned};
    \node[font=\scriptsize, below=2mm of scorers.south] {policy-owned};
  \end{tikzpicture}

  \vspace{4pt}
  \begin{minipage}{0.7\linewidth}
\begin{footnotesize}
\begin{verbatim}
worker_selection:
  aggregated: candidate      # instance for full requests
  instances:
    - name: candidate
      type: ramjet           # any catalog type
      parameters:
        alpha: 0.883
        max_affinity_blocks: 68
        load_unit_tokens: 4552
\end{verbatim}
\end{footnotesize}
  \end{minipage}
  \caption{Top: one worker selection. The host fixes the eligible candidates and fills the inputs the
    policy's components declared; the policy's filters, scorers and picker turn them into one row.
    Bottom: the policy file that selects a catalog type and its parameters; the parameters shown are
    the tuned \policy{ramjet} configuration that serves as the best baseline in this study.}
  \label{fig:api}
\end{figure}
% src (ramjet parameters): CR/facts/finalists.json headline.val_best_baseline.spec.parameters
%      (alpha 0.8828774076622825, max_affinity_blocks 68, load_unit_tokens 4552,
%      affinity_block_tokens 512, basis relative)

% ---------------------------------------------------------------------------------------------
\subsection{Dynamo's default cost function}
\label{sec:default-cost}

% src: WT/lib/router-plugins/builtin/src/default/scorer.rs (score_worker, aggregated branch);
%      knobs in default/parameters.rs. The formula below is the aggregated, prefill-tracking path
%      replay uses; decode-only, host/disk-tier, shared-cache and taint terms are omitted.
In the API's terms, the default router (\defaultcf{} in the catalog) is one scorer and one picker.
The scorer declares \code{LOAD}, \code{PREFERRED\_TAINT} and, unless its overlap credit is zero,
\code{CACHE}, and assigns each candidate a cost in KV blocks:
\begin{equation}
  \label{eq:default-cost}
  \cdef_i = \bpf \max\!\Big(0,\ \frac{\Lpf_i + \ISL_j}{16} - \bov\,\gamma_i\,\ovl_i\Big)
           + \Ldec_i + \breq \Lreq_i,
\end{equation}
where $\ovl_i$ is the candidate's device-tier overlap in blocks, $\Lpf_i$ its active prefill tokens,
$\Ldec_i$ its decode cost blocks (including this request's new blocks), $\Lreq_i$ its active
requests, 16 the KV block size, and
\(
  \gamma_i = \big(1 + \bdecay\,(\Lpf_i - \min_{i' \in \cand}\Lpf_{i'})/(16\,\reqblocks_j)\big)^{-1}
\)
discounts the overlap credit of workers with more queued prefill than the least-loaded one. The
shipped defaults are an overlap credit $\bov = 1$, a decay $\bdecay = 0$, a prefill load scale
$\bpf = 1$ and a decode request weight $\breq = 0$, so the default router minimizes the prefill
blocks the worker would still have to compute plus its decode load.\footnote{The full scorer also
  credits host- and disk-tier and shared-cache hits, has a separate branch for decode-only workers,
  and multiplies the cost by a preferred-taint factor; none of these enter the aggregated,
  device-tier-only setting studied here, in which replay requests carry no taint preferences.}
% src (taint): WT/lib/mocker/src/replay/offline/extensions/kv_router/mod.rs (routing_constraints:
%      RoutingConstraints::default(), so preferred_taints is empty and the multiplier is absent)

The picker takes the minimum cost; with a \code{router\_temperature} above zero it samples from a
softmax over costs normalized to their range across the candidates, so at positive temperature the
default is set-dependent. \Cref{app:api} gives its tie-breaking, seeding and session-affinity
details.
% src: default/picker.rs (softmax_sample_index: range-normalized costs)

% ---------------------------------------------------------------------------------------------
\subsection{The ported cost functions as instances of the API}
\label{sec:baseline-policies}

\Cref{tab:api-policies} writes every catalog policy type of the study in the API's vocabulary: its
components, the inputs and request facts it reads, whether its choice for one worker depends on the
other candidates (set dependence), the state it keeps, and the parameters this study tunes. The
pedagogical value of the table is that the policies differ along a few axes only. Most are a single
picker over \code{CACHE} and \code{LOAD}; they differ in how they combine overlap with load
(additive, multiplicative, thresholded or lexicographic), in whether they restrict candidates by a
hash of the prompt prefix, and in whether they compare a worker with the rest of the set. The
router's round-robin mode is not a catalog policy: it sends requests to the workers in turn, reads
neither cache nor load, has no parameters, and serves as a floor.
% src: CR/CONTRACT.md "Policy YAML" (round-robin uses router_mode="round_robin", no YAML)
\Cref{app:ports} states each rule in full and lists where the ports differ from the published
methods; \cref{sec:ports} discusses the difference that bears on the headline.

% src: per-row components and inputs from the WorkerSelectionPolicy::new calls and the
%      required_worker_inputs of WT/lib/router-plugins/builtin/src/{default.rs, default/scorer.rs,
%      two_tier_cost_fn.rs, lmetric.rs, ramjet.rs, dualmap.rs, chwbl.rs, llm_d/precise_prefix.rs,
%      llm_d/optimized_baseline.rs, sticky_session.rs, thunderagent/worker_selection/mod.rs,
%      learned_choice/mod.rs}; request-fact accessors grepped per file (prefix_hashes: chwbl,
%      dualmap, lmetric, learned_choice v2; session_context: sticky_session, learned_choice;
%      policy_class and worker_capacity: precise_prefix; worker_capacity: learned_choice;
%      affinity_target: thunderagent); exclusive affinity: default, sticky_session, learned_choice;
%      tuned parameters: CR/facts/phase2_prep.json spaces (tab:baseline-spaces)
{\footnotesize
\setlength{\LTcapwidth}{\linewidth}
\setlength{\tabcolsep}{4pt}
\begin{longtable}{@{}P{0.17\linewidth}P{0.18\linewidth}P{0.14\linewidth}P{0.18\linewidth}P{0.12\linewidth}P{0.1\linewidth}@{}}
\caption{Catalog policy types as compositions of API components. S: scorer; P: picker (every
  policy has exactly one). Inputs: declared \code{WorkerInputs} (C = \code{CACHE}, L = \code{LOAD},
  PT = \code{PREFILL\_TIME}, T = \code{PREFERRED\_TAINT}). Request facts beyond the prompt length:
  hashes = prefix hashes, session = session context, capacity = advertised KV capacity.
  Set-dependent: whether a candidate's score or eligibility depends on the other candidates at the
  shipped parameters.}
\label{tab:api-policies}\\
\toprule
Type (variants) & Components & Inputs, request facts & Set dependence & State & Tuned here \\
\midrule
\endfirsthead
\toprule
Type (variants) & Components & Inputs, request facts & Set dependence & State & Tuned here \\
\midrule
\endhead
\bottomrule
\endfoot
\code{dynamo-default-}\allowbreak\code{cost-fn} & S default cost; P min cost (softmax if $\temp > 0$) & C, L, T &
  through $\gamma_i$ if $\bdecay > 0$, and the softmax & random generator & 5 knobs (M0,
  \cref{sec:ladder}) \\
\code{dynamo-two-tier-}\allowbreak\code{cost-fn} & P lexicographic (load tier, cache tier) & C, L &
  yes: request spread and best overlap across the set & none & 3 thresholds \\
\policy{lmetric} & P min of a product & C, L, hashes &
  hot-spot classes (holders against non-holders) & class window & 2 \\
\policy{ramjet} & P max of affinity minus load & C, L &
  yes: affinity relative to the warmest candidate & tie rotation & 3 \\
\policy{dualmap} & P two hashed candidates, warmer unless over budget & C, L, hashes &
  yes: candidates fixed by hash & prefix window & 3 \\
\policy{chwbl} & P hashed owner within a load bound & L, hashes &
  yes: bound relative to the set mean & none & 2 \\
\code{llm-d-precise-}\allowbreak\code{prefix} & S prefix, S queue, S KV utilization; P max score & C, L,
  capacity, policy class & yes: queue score scaled by the set's range & tie rotation & 4 \\
\code{llm-d-optimized-}\allowbreak\code{baseline} (throughput, modeled) & S token load; P sticky filter and TTFT gate &
  C, L (PT if modeled) & yes: sticky set and TTFT gate & tie rotation & 3 \\
\policy{sticky-session} (hard, bounded) & S default cost; P session's worker or min cost &
  C, L, T, session & bounded: request count against the set mean & session map, generator &
  5 or 6 \\
\policy{thunderagent} & S prefill and decode load; P classifier's target or min cost & L, affinity
  target & no & in its request classifier & not evaluated \\
\midrule
\lc{} (this work) & S default cost (feature 0); P utility argmax & C, L, T, capacity, session,
  hashes (\code{v2}) & \code{v1} at rank 0: no; \code{v2} and context terms: yes & session map,
  generator & 7 to 22 coefficients (headline league) \\
\end{longtable}
}

% ---------------------------------------------------------------------------------------------
\subsection{The learned policy as a generalization}
\label{sec:api-learned}

Each of these rules hand-picks a functional form over a few of the declared inputs. \lc{}
(\cref{sec:model}) instead scores every candidate with one utility that is linear in a fixed vector
of features computed from the same inputs, and lets the coefficients be tuned. It is built from the
same parts as the rules it generalizes: the default cost function's own scorer supplies its first
feature, and a picker computes the rest and takes the argmax. Several of the heuristics are points,
or limits, of its parameter space.
\begin{itemize}
  \item $\theta = -\basis{0}$ is the default router (\cref{sec:parity}).
  \item A large positive session coefficient is hard stickiness (\cref{sec:sticky}).
  \item With the extended feature set \code{v2}, $\theta = -(\basis{8} + \basis{10})$ takes the
    argmin of LMetric's product with the queued-prefill term, and $-(\basis{9} + \basis{10})$
    reproduces the branch's \code{lmetric} port as evaluated, without its hot-spot filter
    (\cref{sec:features}).
  \item \code{v2}'s \code{hash\_home} indicator marks the two rendezvous-hashed candidates of a
    prompt prefix, the candidate set of DualMap and Lodestar \citep{yuan2026dualmap,lim2026lodestar}.
  \item Set-relative features and context terms let the comparison between two workers depend on
    the others (\cref{sec:context}), as the queue score of \policy{llm-d-precise-prefix} or the
    relative basis of \policy{ramjet} do by construction.
\end{itemize}
What it cannot express are hard thresholds and lexicographic rules (two-tier's tiers, DualMap's
budget, CHWBL's bound, the optimized baseline's TTFT gate) and history-dependent windows, which a
smooth utility can only approximate. The study's question is whether tuning a generalization
of this kind, with the same budget as each hand-picked rule, finds a better trade on one deployment.

% ---------------------------------------------------------------------------------------------
\subsection{Caveats that bound the results}
\label{sec:caveats}

% src: rows 1, 2, 3, 9, 10, 13 and 15 of tab:caveats (sections/appendix-caveats.tex), whose row
%      comments name the evidence
Building, calibrating and auditing the study surfaced properties of the API, the replay host and the
workloads that can change which policy wins, and several of them shape the results. Replay's router
state is perfectly fresh, while a live router sees cache events and load late (caveat 1;
\cref{sec:res-robustness,sec:live}). Replay fills \code{expected\_output\_tokens} with the true output
length, a value no live router has; no learned feature reads it (caveat 2). Host knobs change what
the inputs mean: without assumed KV reuse every prefix hash is unique to its request, so prefix-keyed
policies lose their keys, and the learned policy must run with its training-time knobs (caveat 3;
\cref{sec:threats-scope}). Logged arrival timestamps are quantized, so the harness spreads arrivals
per replicate (caveat 9; \cref{sec:replicates}). No AgentX play (one recorded agent run;
\cref{sec:workloads}) fits a 32K context, so the deployment runs at 128K with YaRN context extension
(a factor of 4 over the model's original \num{32768} tokens), where \ais{}'s accuracy is least
validated, and \aisim{} prices a mixed prefill
pass at the batch mean, an artifact a search could exploit (caveats 10 and 13;
\cref{sec:threats-sim,sec:gaming}). Finally, the request-count objective can be gamed (caveat 15;
\cref{sec:gaming}). \Cref{app:caveats} lists all 15 with their evidence and status (fixed, follow-up,
inherent).
% src (YaRN gloss): CR/facts/live.json ("--max-model-len 131072 with YaRN (factor 4, original 32768)");
%      CR/config/engine.json max_model_len note ("131072 per operator (YaRN)"); restructure round 1, E4
% src (mixed prefill pricing): CR/facts/UPSTREAM_FOLLOWUPS.md #17
%      (AISim, the replay engine, prices a mixed prefill pass; it calls AIS predict_prefill on batch means)
